%% kapa: exam=1 round=36 year=2025 period=1 form=A subject=civil_law subject_ko=민법 seq=9 number=9
%% source: Q-Net 감정평가사 기출문제 (cst00309) artlSeq=5231525 file=raw/1008/5231525/2025년 제36회 감정평가사 1차시험 1교시 문제지 원본.pdf
%% conversion: kordoc 4.24.51 -> tools/gen1.py (auto); review=proofread-sonnet; flags=none
\begin{question}{9}
물건에 관한 설명으로 옳지 않은 것은? (다툼이 있으면 판례에 따름)
\begin{choices}
  \item 기존 건물에 증축을 한 소유자의 구분행위가 없더라도 증축 부분이 구조상ㆍ이용상의 독립성을 갖춘 경우, 그 부분에 대한 구분소유권은 성립된다.
  \item 건물의 개수는 사회통념 또는 거래관념에 따라 물리적 구조 등 객관적 사정과 건축한 자의 의사 등 주관적 사정을 참작하여 결정된다.
  \item 타인 소유의 물건은 종물이 될 수 없다.
  \item 당사자는 주물을 처분할 때에 특약으로 종물을 제외할 수 있다.
  \item 매매목적물의 인도 전에도 매수인이 매매대금을 완납한 후 발생한 그 목적물의 과실에 대한 수취권은 특별한 사정이 없는 한 매수인에게 귀속된다.
\end{choices}
\end{question}
\answer{1}
